\section{Which Opportunity Deserves the First Experiment?}

\subsection{1st Pass: Non-Negotiable Hard Gates}

Once you have a handful of opportunity cards, comparison becomes unavoidable. There is no universal mathematical formula. Markets are uncertain, and assessments are qualitative. Still, disciplined evaluation beats chasing trends.

Before scoring, eliminate ideas that fail a non-negotiable gate:

\begin{itemize}[leftmargin=11pt,itemsep=1.5pt,topsep=1.5pt,parsep=0pt]
  \item The proposed activity cannot be conducted legally or ethically in your current circumstances.
  \item The 1st meaningful experiment exceeds your affordable-loss budget.
  \item The business model requires daytime responsiveness you cannot supply.
  \item Delivery depends on proprietary, confidential, or non-permitted employer resources.
  \item You cannot identify any plausible route to contacting target economic buyers.
\end{itemize}

A failed gate means \textit{defer and learn}, not \textit{never}. But a low score on a gate cannot be compensated by high scores elsewhere.

\subsection{2nd Pass: Compare Uncertainties, Not Fantasies}

\textbf{The most critical distinction is between what is attractive about a market in the abstract and what makes it a plausible opening for you right now.}

\begin{tcolorbox}[
  enhanced,
  colback=white,
  colframe=BayesBorder!80!white,
  boxrule=0.4pt,
  arc=2pt,
  left=3pt,right=3pt,top=3pt,bottom=3pt,
  before skip=3pt,after skip=5pt
]
\scriptsize
\setlength{\tabcolsep}{2pt}
\renewcommand{\arraystretch}{1.18}
\begin{tabularx}{\linewidth}{@{} >{\bfseries\color{BayesInk}}p{1.85cm} >{\color{BayesInk}}X >{\color{BayesSecondary}}p{2.3cm} @{}}
\toprule
\textbf{Dimension} & \textbf{Evaluation Question} & \textbf{What Counts as Proof?} \\
\midrule
Problem severity & Is this costly enough to buy? & Observed spend; workarounds. \\
Founder insight & Do I see hidden trade-offs? & Correct diagnostic verdicts. \\
Feasibility & Can I build it rapidly? & Completed work samples. \\
Buyer access & Can I reach decision-makers? & Verified meetings, intros. \\
Credibility & Can I reduce perceived risk? & Unseen demos, references. \\
Time-to-learn & Can I test in 2 weeks? & Narrow diagnostic trial. \\
Economics & Healthy contribution margin? & Bottom-up time/cost audit. \\
Motivation & Will I endure 6 hard weeks? & Past resilience in domain. \\
\bottomrule
\end{tabularx}
\end{tcolorbox}

\subsection{Comparison Without Fake Precision}

Replace decimal scores with qualitative evidence tags: \textbf{Observed}, \textbf{Corroborated}, \textbf{Speculative}, and \textbf{Blocked}.

\begin{tcolorbox}[
  enhanced,
  colback=white,
  colframe=BayesBorder!80!white,
  boxrule=0.4pt,
  arc=2pt,
  left=3pt,right=3pt,top=3pt,bottom=3pt,
  before skip=3pt,after skip=5pt
]
\scriptsize
\setlength{\tabcolsep}{2pt}
\renewcommand{\arraystretch}{1.18}
\begin{tabularx}{\linewidth}{@{} >{\bfseries\color{BayesInk}}p{1.7cm} *{3}{c} >{\color{BayesSecondary}}X @{}}
\toprule
\textbf{Candidate} & \textbf{Demand} & \textbf{Fit} & \textbf{Access} & \textbf{Next Experiment} \\
\midrule
AI Eval Service & Spec & Corr & Spec & Interview 8 budget owners \\
Dev Cohort & Spec & Corr & Spec & Paid workshop + exam \\
Data Center & Unclear & Weak & Blocked & Shelve; capital barrier \\
\bottomrule
\end{tabularx}
\end{tcolorbox}

\subsection{The Kill-Question Method}

For each candidate, formulate 3 falsification questions:
\begin{enumerate}[label=\textbf{\arabic*.},leftmargin=11pt,itemsep=1.5pt,topsep=1.5pt,parsep=0pt]
  \item Do buyers already allocate a budget or urgent time to this pain?
  \item Can I produce a verifiable improvement within my affordable loss?
  \item Can I reach enough buyers to learn whether demand is repeatable?
\end{enumerate}

\textbf{Use the kill-question method: prioritize experiments that could decisively falsify the premise, rather than busying yourself with logos and incorporation.}
